\section{Reduced costs and directions from subtree sums}\label{sec:directions}

We now express one revised simplex iteration using the forest. The positive
tree is $T_*$, its weight is $W_*=w(T_*)$, and the current objective is
$\rho=p(T_*)/W_*$. All formulas refer to the current basis.

\subsection{Connection with revised simplex: the two basis systems}
Write $\mat{B}=\mat{A}_{\Bset}$ and let $\vect{g}=(\vect{p},0)$ be the standard-form objective vector.
The two familiar basis systems in a primal revised simplex iteration are
\begin{align}
 \mat{B}^T\vecg{\pi}&=\vect{g}_{\Bset},
 &\bar c_q&=g_q-\mat{A}_q^T\vecg{\pi},
 &&\text{multipliers and reduced costs},\label{eq:pricing-system}\\
 Bv&=A_q,
 &d_{\Bset}&=-v,
 &&\text{column for entering variable }q.\label{eq:column-system}
\end{align}
The forest supplies explicit solutions to both systems. For the first,
subtree sums yield the reduced costs in \eqref{eq:rc-node}--\eqref{eq:rc-edge}
and the multipliers $\pi=(\rho,\alpha)$ in \eqref{eq:forest-dual}. For the
second, a signed tree or subtree movement and one normalization correction
give \eqref{eq:direction}. Thus neither stage requires a general basis
factorization or a triangular solve. The next subsections derive these
identities and explain their arithmetic costs. In this precise sense, Forest
Simplex implements the revised simplex operations through a combinatorial
representation of the basis: it computes the same multipliers, reduced costs,
and pivot directions by forest traversals and aggregate formulas.

The entering variable selects a \emph{column}; the ratio test subsequently
selects the leaving variable and hence a \emph{row}. If that pivot row itself
is required, its usual transposed basis system also has a direct forest
evaluation, given in \eqref{eq:pivot-row} below.

\subsection{The dictionary behind the formulas}
Let $a_T$ be the root variable of tree $T$. In zero trees this is the
nonbasic root $y_{q_T}$; in the positive tree it is basic. Repeated use of
\eqref{eq:parent-child} gives
\begin{equation}\label{eq:path-dictionary}
 y_i=a_{T(i)}+\sum_{e\text{ on root-to-}i\text{ path}}\sigma_e s_e.
\end{equation}
Summing with weights, each $s_e$ appears exactly at the descendants $D_e$.
Therefore
\begin{equation}\label{eq:root-dictionary}
 a_{T_*}=\frac{1-\sum_{T\ne T_*}w(T)y_{q_T}
                    -\sum_{e\in\Fset}\sigma_e w(D_e)s_e}{W_*}.
\end{equation}
Equations~\eqref{eq:path-dictionary}--\eqref{eq:root-dictionary}, followed
by $s_{ij}=y_j-y_i$ for nonforest arcs, are an explicit simplex dictionary.
They express every basic variable in terms of the $n-1$ nonbasic variables.
Only the positive root is eliminated by normalization.

Substituting into the objective yields
\begin{equation}\label{eq:objective-dictionary}
\begin{split}
 \vect{p}^T\vect{y}={}&\rho+
 \sum_{T\ne T_*}\bigl[p(T)-\rho w(T)\bigr]y_{q_T}\\
 &+\sum_{e\in\Fset}\sigma_e\bigl[p(D_e)-\rho w(D_e)\bigr]s_e.
\end{split}
\end{equation}
This identity immediately gives all reduced costs, without solving a dense
linear system.

\begin{proposition}[Forest pricing]\label{prop:pricing}
For the maximization problem, the reduced costs of the nonbasic variables are
\begin{align}
 \bar c_{y_{q_T}}&=p(T)-\rho w(T),&&T\ne T_*,\label{eq:rc-node}\\
 \bar c_{s_e}&=\sigma_e\bigl[p(D_e)-\rho w(D_e)\bigr],&&e\in\Fset.
 \label{eq:rc-edge}
\end{align}
A positive reduced cost identifies a candidate entering variable. It does
not guarantee a positive feasible step.
\end{proposition}

\paragraph{Specialized entering-variable selection.}
A postorder traversal computes all subtree profit and weight sums in $O(n)$
arithmetic operations. With these aggregates available, either formula above
costs $O(1)$ per candidate. There are exactly $n-1$ nonbasic variables, so full
pricing and Bland's entering-variable selection cost $O(n)$, independently of
the number m of original arcs. No arc outside the forest needs pricing.
Writing $P_*=p(T_*)$, signs and comparisons of reduced costs can also use
their numerators over the common positive denominator $W_*$:
\[
 W_*p(T)-P_*w(T),\qquad
 \sigma_e\bigl[W_*p(D_e)-P_*w(D_e)\bigr].
\]
These comparisons avoid division; they do not remove the need to control
overflow or rounding in a numerical implementation.

\subsection{A specialized pivot column from one component movement}
Choose a nonbasic variable $q$ and increase it from zero to $t$. Define the
unadjusted vertex displacement $h$ by
\begin{equation}\label{eq:h}
 h=\begin{cases}
 \ind^T,&q=y_{q_T}\text{ is a root},\\
 \sigma_e\ind^{D_e},&q=s_e\text{ is a forest slack}.
 \end{cases}
\end{equation}
The raw motion generally changes total weight. Restore normalization by
shifting every vertex of $T_*$ by the same amount:
\begin{equation}\label{eq:direction}
 \kappa=\frac{\sum_i w_i h_i}{W_*},\qquad
 d_i=h_i-\kappa\ind^{T_*}_i,\qquad
 d_{s_{ij}}=d_j-d_i\quad((i,j)\in\Aset).
\end{equation}
The full candidate motion is
$y_i(t)=y_i+t d_i$ and $s_{ij}(t)=s_{ij}+t d_{s_{ij}}$.

\begin{proposition}[The forest motion is the simplex direction]
\label{prop:direction}
The direction in \eqref{eq:direction} has component one at the entering
variable, zero at every other nonbasic variable, and satisfies all homogeneous
equations of \eqref{eq:standard}. Consequently, its basic components equal
$-A_{\Bset}^{-1}A_q$, and its objective derivative is $\bar c_q$.
\end{proposition}
\begin{proof}
For a root entry, $h$ is constant on each tree, and the chosen root
increases by one. All other roots remain fixed. For an edge entry, $h$
is constant on each side of the deleted tree edge and its head-minus-tail
difference on that edge is one by \eqref{eq:sign}. No other forest edge has
a nonzero head-minus-tail difference. No zero-tree root belongs to $D_e$.
The correction on $T_*$ is constant on that entire tree and affects neither
these differences nor any zero-tree root. By construction $\vect{w}^T\vect{d}=0$, and
the slack differences enforce every arc equation. Uniqueness follows from
nonsingularity of the basis. Finally,
\[
 \vect{p}^T\vect{d}=\vect{p}^T\vect{h}-\rho \vect{w}^T\vect{h},
\]
which gives \eqref{eq:rc-node} or \eqref{eq:rc-edge}.
\end{proof}

\paragraph{Computing the column without a basis solve.}
The conventional tableau column is $A_{\Bset}^{-1}A_q=-d_{\Bset}$.
Equation~\eqref{eq:direction} computes it directly from the forest, including
its sign convention. For root entry, $\kappa=w(T)/W_*$; for slack entry,
$\kappa=\sigma_e w(D_e)/W_*$. Each uses a stored aggregate.
The direction is described by the selected tree or subtree, its sign, the
positive tree, and one scalar. Once component membership is available and
the selected subtree is marked, each vertex coordinate and each slack
coordinate costs $O(1)$ to evaluate. Marking and forming all vertex coordinates
cost $O(n)$; explicitly forming all slack coordinates costs $O(m)$.
The ratio test can evaluate these coordinates as it scans the basic variables,
so storing a separate complete pivot column is optional.

\paragraph{The pivot row, explicitly.}
Suppose the ratio test chooses basic variable $z_b$, at position r in the
ordered basis. A conventional way to obtain its tableau row is to solve
$\mat{B}^T\vecg{\eta}=\vect{e}_r$ and form $a_{rq}=\vecg{\eta}^T\mat{A}_q$.
For each nonbasic candidate q, let $h^{(q)}$ and $\kappa_q$ be the quantities
in \eqref{eq:h}--\eqref{eq:direction}. The same forest formulas give directly
\begin{equation}\label{eq:pivot-row}
 a_{rq}=-d_b^{(q)}=
 \begin{cases}
 \kappa_q\ind^{T_*}_i-h_i^{(q)},&z_b=y_i,\\[2pt]
 h_i^{(q)}-h_j^{(q)}+
 \kappa_q(\ind^{T_*}_j-\ind^{T_*}_i),&z_b=s_{ij}.
 \end{cases}
\end{equation}
Indeed, $d_{\Bset}^{(q)}=-B^{-1}A_q$, so its component at position r
is the negative of the requested row coefficient. Component identifiers and
depth-first entry/exit indices allow tree and subtree membership queries in
$O(1)$ after $O(n)$ forest preprocessing. With the aggregates also available,
each coefficient above costs $O(1)$ and all $n-1$ nonbasic coefficients of a
chosen row cost $O(n)$. The basic coefficients are the corresponding identity
row. The reference primal algorithm only needs the entering column and ratio
test; \eqref{eq:pivot-row} makes explicit how to recover the row when needed.

\paragraph{Why all arcs still matter.}
Only forest arcs are needed for pricing and vertex directions. However, the
basic slack on \emph{every} nonforest arc can decrease. Those arcs must be
considered in the ratio test even when they join two different trees or are
transitively redundant. Ignoring them can produce an infeasible pivot.

\subsection{Reading a dual certificate from the forest}
Define arc numbers and vertex residuals by
\begin{align}
 \beta&=\rho,\qquad
 \alpha_e=\begin{cases}-\bar c_{s_e},&e\in\Fset,\\0,&e\notin\Fset,
 \end{cases}\label{eq:forest-dual}\\
 r_i&=w_i\rho+\divg_i(\alpha)-p_i.\label{eq:dual-residual}
\end{align}

\begin{proposition}[Optimality certificate]\label{prop:certificate}
For a forest basis, $r_i=0$ at every basic vertex variable and
$r_{q_T}=-\bar c_{y_{q_T}}$ at every root. Thus if all nonbasic reduced
costs are nonpositive, \eqref{eq:forest-dual} is feasible for
\eqref{eq:ratio-dual} and certifies the optimal ratio $\rho$.
\end{proposition}
\begin{proof}
Set $u_i=p_i-\rho w_i$. On any rooted subtree $D_e$, the only forest edge
crossing its boundary is $e$. Its net outgoing flow is
$-\sigma_e\alpha_e=u(D_e)$ by \eqref{eq:rc-edge}. Taking differences between
a subtree and its child subtrees proves $\divg_i(\alpha)=u_i$ at every
nonroot vertex. On $T_*$, $u(T_*)=0$, so the same equality holds at its root.
In a zero tree, all nonroot equalities hold and the residual at the root
is $-u(T)=\rho w(T)-p(T)$. Nonpositive reduced costs make every $\alpha_e$
and every $r_i$ nonnegative. The primal and dual objectives both equal
$\rho$, proving optimality.
\end{proof}
This certificate uses sums on the original arcs of the residual instance.
It is suitable for a verifier independent of the data structures used to
update the forest.
At every forest basis, even before optimality, the vanishing residuals at
basic vertex variables and the zero multipliers on basic slacks imply
$\mat{B}^T\vecg{\pi}=\vect{g}_{\Bset}$. Thus the construction also solves the first basis system
\eqref{eq:pricing-system}; nonpositive reduced costs are needed only to make
these multipliers dual feasible.
